\chapter{The BGFX Backend}
\label{ch:bgfx-backend}

\texttt{BGFX} integrates the third-party \texttt{bkaradzic/bgfx} rendering library (fetched
directly via CMake's \texttt{FetchContent}, Chapter~\ref{ch:backend-architecture}) behind the
same contract every other backend implements, using bgfx's own native API --- window and
platform initialization, texture creation, sprite draws, frame submission --- rather than
routing through SDL\_Renderer.

\section{Test baseline and the Phase 72 closure}

At the time of writing, the full unit suite passes 4{,}375 of 4{,}377 tests (two hardware
skips), and the backend-specific \texttt{ctest} suite passes 103 of 105. The two remaining
failures are both environment ceilings, not code defects: one render-target MSAA-resolve test
fails because this project's own headless \texttt{Xvfb} display has no DRI3 support to
resolve against, and one render-target-cube depth-format test (a
\texttt{Depth24Stencil8}-attached cube face producing no color output) has been investigated
three separate times without finding a root cause, and remains formally deferred rather than
silently dropped. A first complete, row-by-row triage of this backend originally found
thirty-eight real gaps; thirty-seven were fixed in a single closing session, leaving only the
depth-bias item covered below --- which was itself subsequently fixed too, closing the phase
in full.

\section{Choosing bgfx's underlying native renderer}
\label{sec:bgfx-renderer-selection}

Unlike every other backend in this book, \texttt{BGFX} is itself an abstraction over several
further, distinct native graphics APIs --- OpenGL, OpenGL ES, Vulkan, Direct3D~11/12, and
Metal --- and CNA exposes a real, previously-undocumented mechanism to choose among them: the
\texttt{CNA\_BGFX\_RENDERER} environment variable, resolved through three small, dedicated
functions in \texttt{BgfxRendererSelection.cpp}. \texttt{GetDefaultRendererType()} picks
\texttt{bgfx::RendererType::OpenGL} on Linux specifically and \texttt{Count} (bgfx's own
``let the library auto-detect'' sentinel) everywhere else; \texttt{ParseRendererTypeOverride()}
case-insensitively accepts \texttt{AUTO}, \texttt{OPENGL}, \texttt{OPENGLES}, \texttt{VULKAN},
\texttt{METAL}, \texttt{DIRECT3D11} / \texttt{D3D11}, \texttt{DIRECT3D12} / \texttt{D3D12}, and
\texttt{NOOP} (bgfx's own headless no-op renderer, useful for pure logic testing with no GPU
work issued at all), throwing a \texttt{std::runtime\_error} naming every valid value for
anything else; \texttt{ResolveRendererType()} ties the two together, falling back to the
platform default whenever the environment variable is unset or empty.

\subsection{A worked example: routing one specific test around a real bgfx/GL limitation}

This mechanism exists for a concrete, real reason, not as a speculative escape hatch:
\texttt{examples/bgfx\_rendertarget2d\_msaa\_test.cpp}'s own header comment documents a genuine
investigation. In this project's own Linux development sandbox, bgfx's default OpenGL renderer
negotiates only a legacy OpenGL~2.1 context --- confirmed via \texttt{glxinfo} that the
underlying Mesa/RadeonSI driver actually supports OpenGL~4.6 with full hardware acceleration,
so this is bgfx itself requesting an old context by default, not a driver ceiling --- and
MSAA-flagged framebuffer-attached textures do not actually resolve with real sub-pixel blending
under that legacy GL~2.1 path, despite \texttt{bgfx::getCaps()} itself reporting
\texttt{BGFX\_CAPS\_FORMAT\_TEXTURE\_FRAMEBUFFER\_MSAA} support for the exact format requested.
The investigation ruled out a CNA-side defect specifically: forcing bgfx onto its Vulkan
renderer instead (\texttt{CNA\_BGFX\_RENDERER=VULKAN}, routing through the identical GPU) made
both of the test's own checks pass cleanly, including a genuine intermediate (blended) pixel
value along the test triangle's diagonal edge, with \emph{zero} code changes on either side ---
proof the C++ wiring itself is correct and backend-agnostic, and the gap is specifically in
bgfx's own default legacy-GL-context resolve path. The real \texttt{ctest} registration
(\texttt{cmake/Tests/BgfxTests.cmake}) reflects this finding directly rather than leaving it as
a comment nobody acts on:

\begin{lstlisting}
cna_register_backend_test(NAME Bgfx_RenderTarget2D_MsaaResolve
    COMMAND cna_test_bgfx_rendertarget2d_msaa
    TIMEOUT 30
    ENVIRONMENT "SDL_VIDEODRIVER=x11;DISPLAY=${CNA_TEST_DISPLAY};CNA_BGFX_RENDERER=VULKAN")
\end{lstlisting}

This is worth reading alongside, not instead of, this chapter's own ``Test baseline'' section
above and \texttt{docs/graphics-backend-feature-matrix.md}'s repeatedly-reconfirmed count
(spanning multiple independently-dated task closures) that \texttt{Bgfx\_RenderTarget2D\_%
MsaaResolve} remains one of exactly two still-failing tests in \emph{that} project's own
tracked baseline, attributed there to a related but distinctly-framed cause: ``this sandbox's
\texttt{Xvfb} has no DRI3 support.'' The two explanations are not actually in tension --- DRI3
is specifically what a windowed OpenGL context needs to negotiate real hardware acceleration
under X11 indirect rendering, so its absence is a plausible reason bgfx's OpenGL path falls
back to a legacy context in the first place, while Vulkan's own presentation integration does
not depend on the same DRI3/GLX negotiation at all. What this book cannot independently confirm
is whether the specific sandbox behind that repeatedly-cited count genuinely has real Vulkan/GPU
access available to the \texttt{CNA\_BGFX\_RENDERER=VULKAN} override the way the test's own
comment describes, or whether that particular environment lacks real GPU access down either
path --- this book's own screenshot-feasibility investigation (Chapter~\ref{ch:graphicsdevice})
independently confirmed no GPU passthrough at all in \emph{this} book-writing session's own
container, a genuinely different environment from CNA's own separate development sandbox. Both
real, sourced explanations are given here rather than silently picking one, consistent with
this book's own methodology of not resolving an ambiguity the underlying sources themselves
leave open.

\section{DepthBias: a real gap in the underlying library, not just CNA}

\cnaclass{RasterizerState.DepthBias} is the one feature in this chapter whose fix required
working around a genuine absence in bgfx itself, confirmed by reading bgfx's own vendored
source directly: its high-level state API has no depth-bias mechanism of any kind --- no
corresponding flag in its own headers, and no \texttt{glPolygonOffset} call anywhere in its
vendored OpenGL renderer backend. The fix emulates the feature entirely at the CNA level, via
a per-draw vertex-shader Z-offset (a new \texttt{u\_depthBias} uniform threaded through every
3D vertex shader). \texttt{SlopeScaleDepthBias}, by contrast, remains \emph{deliberately}
unimplemented --- a project-owner decision, not an oversight --- because a true
per-fragment, screen-space-slope computation would force every 3D shader off the GPU's
early-Z optimization path, a cost judged not worth paying for this specific, rarely-used
parameter.

\section{Render-target readback crashes: one root cause, five fixes}

Five of six render-target-related \texttt{glReadPixels} / \texttt{Xvfb} crashes were traced to
a single, shared root cause: bgfx processes its internal views in ascending ID order every
frame, which means any render-target view is always the \emph{last} one processed --- and
therefore still GL-bound at the exact moment \texttt{glReadPixels()} fires for a screenshot.
The fix adds a dedicated ``flush'' view, touched immediately before every screenshot
specifically to force the correct view to be current first. Only the render-target-cube
depth-format crash mentioned above has a different, still-unidentified root cause and was
not resolved by this fix.

\section{Sampling a render target: a wrong-handle-type cast, since fixed}
\label{sec:bgfx-wrong-handle-cast}

Chapter~\ref{ch:vulkan-backend} named this backend's own, now-resolved
\cnaclass{RenderTargetCube}-sampling bug as ``a genuine wrong-handle-cast, diagnosed precisely''
in contrast to Vulkan's own still-unresolved equivalent --- this section is that precise
diagnosis, including how it was actually closed. Both \texttt{BgfxSpriteBatchBackend::Draw} and
\cnaclass{BgfxGraphicsBackend}'s own \texttt{envMapping} branch (the code path
\cnaclass{EnvironmentMapEffect}, per Chapter~\ref{ch:stock-effects}, reaches) used to
unconditionally \texttt{static\_cast} whatever \cnaclass{ITextureBackend} / \cnaclass{%
ITextureCubeBackend} they were handed to the plain, ordinary-texture concrete type.
\cnaclass{RenderTarget2D} and \cnaclass{RenderTargetCube} are each backed by their own,
entirely separate sibling backend class --- consistent with
Chapter~\ref{ch:backend-architecture}'s own \cnaclass{IRenderTargetBackend} interface design,
not a shortcut specific to this backend. Both of those sibling classes'
\emph{first data member} is a framebuffer handle, not a texture handle.

The reason this compiled cleanly and never crashed, rather than failing loudly the moment it was
exercised, is a genuinely instructive piece of type-confusion: bgfx represents both handle kinds
identically at the ABI level, as a bare \texttt{struct \{ uint16\_t idx; \}}. The cast was
structurally valid C\texttt{++} --- both types really do have that exact layout --- so the
compiler had no basis to reject it, and the resulting object did not crash on access, because
\texttt{idx} really was a valid, in-range integer. It was simply the \emph{wrong} integer: an
index into bgfx's framebuffer handle pool, read and used as though it indexed the texture handle
pool instead. The practical consequence was a render target that silently sampled whatever
texture happened to occupy the coincidentally-matching slot in the wrong pool --- confirmed by
direct memory-layout analysis and by new smoke tests added specifically to prove the cast
\emph{did not crash} on either \cnaclass{RenderTarget2D} or \cnaclass{RenderTargetCube}, which is
exactly why ordinary ``does it crash'' testing had not caught this before. Both call sites are
now fixed (Tasks~873/874, closed by Task~907's own follow-up): \texttt{BgfxSpriteBatchBackend::%
Draw} now reaches a render target's real dimensions through \cnaclass{ITextureBackend}'s own
virtual \texttt{GetWidth()} / \texttt{GetHeight()} rather than an unsafe cast, and the
\texttt{envMapping} branch resolves a bound cube texture through a real \texttt{dynamic\_cast} to
a shared \texttt{IBgfxCubeSamplable} interface instead --- a render-to-texture technique (a
portal, a security-camera view, a reflection probe rendered as an ordinary 2D texture rather than
through the dedicated cube-face path) that samples its own render target back through
\cnaclass{SpriteBatch} or \cnaclass{EnvironmentMapEffect} on this backend can now be trusted,
where it previously could not have been even though nothing about running the code signaled that
anything was wrong. This chapter's own closing summary below already reflected this fix
implicitly, by not naming Tasks~873/874 among the backend's remaining open items --- the
paragraph above exists so that omission reads as a confirmed, explained resolution rather than
as an unexplained gap between two sections of the same chapter.

\section{A real, found-and-fixed rendering bug: silent black meshes}

One bug is worth naming for how invisible it was until specifically tested: this backend's
graphics-backend class never overrode the effect-aware indexed-draw entry point at all. Any
indexed draw bound to an \cnaclass{Effect} whose vertex format lacked a \texttt{Color}
attribute --- which describes essentially any \cnaclass{Model} loaded through
\texttt{ContentManager} (Chapter~\ref{ch:models-meshes}) --- silently fell back to the base
class's default implementation, discarding the entire \cnaclass{GpuDrawParams} struct and
reading an unbound color attribute (which GL defaults to black). The practical, visible
consequence: any loaded 3D model rendered as a solid black silhouette on this backend,
regardless of its real diffuse color, texture, or lighting --- fixed by adding a real
override.

\subsection{A worked example: recognizing the trigger condition before hitting it}

The black-mesh bug's real trigger condition --- a vertex format with no \texttt{Color}
attribute, which is essentially every \cnaclass{Model} loaded from content rather than
hand-authored with vertex colors baked in --- is checkable directly against a real
\cnaclass{VertexDeclaration}, grounded in \texttt{VertexElement}'s real accessor:

\begin{lstlisting}
bool HasColorAttribute(const VertexDeclaration& decl)
{
    for (const VertexElement& element : decl.GetVertexElements()) {
        if (element.getVertexElementUsageProperty() == VertexElementUsage::Color) {
            return true;
        }
    }
    return false;
}
\end{lstlisting}

Before this bug's fix, any effect-aware indexed draw where \texttt{HasColorAttribute} would
have returned \texttt{false} was exactly the condition that silently rendered solid black on
this backend --- which, per the finding above, covers essentially every ordinary
\cnaclass{Model} draw, since a typical imported mesh carries position/normal/texture-coordinate
data but no per-vertex color. This is precisely why the bug was invisible for so long despite
being this common: nothing about it depended on an unusual asset or an edge-case vertex
layout, only on whether anyone had specifically pixel-tested a textured, unlit 3D model on
this one backend.

\section[OcclusionQuery: a real implementation, correctness not provable here]{OcclusionQuery: a real implementation whose correctness cannot be fully proven here}

\cnaclass{OcclusionQuery} was previously a pure no-op on this backend; the fix wires the
query handle into bgfx's own dedicated submit-with-occlusion-query overload, across all
twelve 3D-draw submission call sites --- no separate correlation/tagging machinery was
needed the way Vulkan's fix required, since bgfx submits synchronously. Two honestly-documented
caveats remain: first, pixel/query correctness genuinely \emph{cannot} be established in this
project's own sandbox, because sabotaging the fix back to a pure no-op produces identical
\texttt{IsComplete()} / \texttt{PixelCount()} output to the fixed version --- the sandbox's own
software Mesa driver returns a non-null result even for a query that was never submitted
anywhere at all, a real ceiling of the software renderer itself, not a CNA defect. Second,
bgfx's own official usage example attaches an occlusion measurement to a separate, dedicated
view specifically so the query is not polluted by other geometry sharing the same depth
buffer; CNA's current fix instead shares whichever view the game's own 3D draw already
targets --- a real, acknowledged architecture gap against true scene-depth query correctness,
not yet attempted.

\section{Where this leaves the backend}

The project's own closing assessment for this backend is unusually clean: ``BGFX's only
remaining code-level gap is [the depth-bias item], already flagged for a project-owner
decision; the other two items are sandbox/environment ceilings\ldots not further BGFX
backend code work.'' Read alongside Chapter~\ref{ch:easygl-backend} and
Chapter~\ref{ch:vulkan-backend}, BGFX is the backend whose remaining open items are most
cleanly separated into ``needs a decision'' and ``needs a different test environment,'' rather
than ``needs more engineering.''
